\subsection{Connection to the Hubble Tension}
  \label{subsec:hubble_tension}

  The Hubble tension arises from the discrepancy between early-time inferences of
  the Hubble constant and late-time local measurements~\cite{Riess2022SH0ES,Planck2020}.
  Comprehensive reviews indicate that no single explanation has yet achieved
  consensus~\cite{Verde2019Review}.

  Within the present framework, this discrepancy reflects the environmental
  sensitivity of late-time kinematic inference.
  Early-time measurements, anchored in the nearly homogeneous and high-density
  early Universe, recover \(H_{\mathrm{global}}\).
  By contrast, late-time distance ladder measurements are performed in a
  structured Universe and may preferentially sample void-dominated regions.

  If local measurements probe environments operating in the saturated regime,
  Eq.~\eqref{eq:hloc} predicts a systematically enhanced inferred value of the
  Hubble constant,
  \begin{equation}
    H_{\mathrm{loc}} > H_{\mathrm{global}}.
  \end{equation}

  The magnitude of the offset depends on the surrounding large-scale environment
  of both observers and sources.
  This naturally explains why the tension manifests primarily in local
  measurements and not in early-time observables.

  The extent to which local underdensities alone can account for the full observed
  tension remains debated~\cite{Shanks2019Void,Kennworthy2019VoidCritique}.
  The present framework provides a distinct and testable mechanism in which the
  effect arises from saturation of the effective transition structure rather than
  from a modification of the background expansion.

  Because the dynamics derives from an effective potential
  \(\Phi_{\mathrm{eff}}\),
  the framework admits a conserved effective energy and a corresponding modified
  virial relation for stationary systems.
  This provides a consistency check ensuring that equilibrium disk galaxies remain
  dynamically stable in the transition regime.
